% GENERATED FILE. Edit canonical lessons, then run npm run generate:books.
% canonical-source-hash: fnv1a64:49ca3d0c445a1f3d
% canonical-lessons: GU-C395-barth, GU-C395-rijharvesan, GU-C395-bhathu, GU-C395-thasvu, GU-C395-besadvu

\chapter{Berths and Reservations}
\label{ch:gu-berths-and-reservations}
\hlchaptermodality{395}

\begin{chapteropening}
\textbf{By the end of this chapter:} \emph{I can find my berth and stow my luggage.}

Complete the last lesson of chapter 395: I can find my berth and stow my luggage.
\end{chapteropening}

\section[\texorpdfstring{\gu{બર્થ} (barth)}{બર્થ (barth)}]{\gu{બર્થ} (barth) — a berth (on a train)}

\label{lesson:GU-C395-barth}



\begin{quote}
Before the new one: say the Gujarati for to slide (something), then the Gujarati for to catch hold of.
\end{quote}

\subsection*{You'll want to know: \gu{બર્થ}}
\textbf{\gu{બર્થ}} — \emph{barth} — \textquotedblleft{}a berth (on a train)\textquotedblright{}.

\begin{quote}
\gu{મારી} \gu{બર્થ} \gu{બારી} \gu{પાસે} \gu{છે}.
\end{quote}

\begin{quote}
\emph{mārī barth bārī pāse che.}
\end{quote}

\begin{quote}
My berth is by the window.
\end{quote}

A feminine noun.

\begin{cousinweb}
\textbf{\gu{બર્થ}} (\emph{barth}) is borrowed from English \emph{berth}, the bed on a night train.
\end{cousinweb}

\begin{grammarlens}[title={What you've built}]
One more word for getting around.
\end{grammarlens}

\subsection*{Your turn}
\begin{itemize}
\raggedright
  \item \emph{Say it:} \emph{barth}
  \item \emph{Say it:} \emph{barth}, once more
  \item \emph{Say it:} \emph{mārī barth bārī pāse che.} — My berth is by the window.
  \item \emph{Recall:} say the Gujarati for to slide (something), then the Gujarati for to catch hold of, then say \emph{barth} again
\end{itemize}

\begin{tcolorbox}[breakable,colback=teal!4,colframe=teal!35!black,title={Before you move on}]
What does \gu{બર્થ} mean? (\textquotedblleft{}A berth (on a train)\textquotedblright{}.) Say it once more.
\end{tcolorbox}
\section[\texorpdfstring{\gu{રિઝર્વેશન} (rijharveśan)}{રિઝર્વેશન (rijharveśan)}]{\gu{રિઝર્વેશન} (rijharveśan) — a reservation (of a seat)}

\label{lesson:GU-C395-rijharvesan}



\begin{quote}
Before the new one: say the Gujarati for to catch hold of, then the Gujarati for a berth (on a train).
\end{quote}

\subsection*{You'll want to know: \gu{રિઝર્વેશન}}
\textbf{\gu{રિઝર્વેશન}} — \emph{rijharveśan} — \textquotedblleft{}a reservation (of a seat)\textquotedblright{}.

\begin{quote}
\gu{અમે} \gu{એક} \gu{મહિના} \gu{પહેલાં} \gu{રિઝર્વેશન} \gu{કર્યું} \gu{હતું}.
\end{quote}

\begin{quote}
\emph{ame ek mahinā pahelā̃ rijharveśan karyũ hatũ.}
\end{quote}

\begin{quote}
We had made the reservation a month before.
\end{quote}

A neuter noun.

\begin{cousinweb}
\textbf{\gu{રિઝર્વેશન}} (\emph{rijharveśan}) is borrowed from English \emph{reservation}, a seat or berth kept in your name.
\end{cousinweb}

\begin{grammarlens}[title={What you've built}]
One more word for getting around.
\end{grammarlens}

\subsection*{Your turn}
\begin{itemize}
\raggedright
  \item \emph{Say it:} \emph{rijharveśan}
  \item \emph{Say it:} \emph{rijharveśan}, once more
  \item \emph{Say it:} \emph{ame ek mahinā pahelā̃ rijharveśan karyũ hatũ.} — We had made the reservation a month before.
  \item \emph{Recall:} say the Gujarati for to catch hold of, then the Gujarati for a berth (on a train), then say \emph{rijharveśan} again
\end{itemize}

\begin{tcolorbox}[breakable,colback=teal!4,colframe=teal!35!black,title={Before you move on}]
What does \gu{રિઝર્વેશન} mean? (\textquotedblleft{}A reservation (of a seat)\textquotedblright{}.) Say it once more.
\end{tcolorbox}
\section[\texorpdfstring{\gu{ભાથું} (bhāthũ)}{ભાથું (bhāthũ)}]{\gu{ભાથું} (bhāthũ) — food packed for a journey}

\label{lesson:GU-C395-bhathu}



\begin{quote}
Before the new one: say the Gujarati for a berth (on a train), then the Gujarati for a reservation (of a seat).
\end{quote}

\subsection*{You'll want to know: \gu{ભાથું}}
\textbf{\gu{ભાથું}} — \emph{bhāthũ} — \textquotedblleft{}food packed for a journey\textquotedblright{}.

\begin{quote}
\gu{માતાએ} \gu{મુસાફરી} \gu{માટે} \gu{થેપલાંનું} \gu{ભાથું} \gu{બાંધ્યું}.
\end{quote}

\begin{quote}
\emph{mātāe musāpharī māṭe theplā̃nũ bhāthũ bā̃dhyũ.}
\end{quote}

\begin{quote}
Mother packed thepla as food for the journey.
\end{quote}

A neuter noun.

\begin{cousinweb}
\textbf{\gu{ભાથું}} (\emph{bhāthũ}) seems to belong with \gu{ભાત}, cooked rice, which CDIAL lists under Sanskrit \emph{bhaktá}, \textquotedblleft{}a meal, food\textquotedblright{} (CDIAL 9331), beside Punjabi \emph{bhattā}, \textquotedblleft{}a travelling allowance\textquotedblright{}; CDIAL gives no form for \gu{ભાથું} itself, so the link is hedged.
\end{cousinweb}

\begin{grammarlens}[title={What you've built}]
One more word for getting around.
\end{grammarlens}

\subsection*{Your turn}
\begin{itemize}
\raggedright
  \item \emph{Say it:} \emph{bhāthũ}
  \item \emph{Say it:} \emph{bhāthũ}, once more
  \item \emph{Say it:} \emph{mātāe musāpharī māṭe theplā̃nũ bhāthũ bā̃dhyũ.} — Mother packed thepla as food for the journey.
  \item \emph{Recall:} say the Gujarati for a berth (on a train), then the Gujarati for a reservation (of a seat), then say \emph{bhāthũ} again
\end{itemize}

\begin{tcolorbox}[breakable,colback=teal!4,colframe=teal!35!black,title={Before you move on}]
What does \gu{ભાથું} mean? (\textquotedblleft{}Food packed for a journey\textquotedblright{}.) Say it once more.
\end{tcolorbox}
\section[\texorpdfstring{\gu{ઠાંસવું} (ṭhā̃svũ)}{ઠાંસવું (ṭhā̃svũ)}]{\gu{ઠાંસવું} (ṭhā̃svũ) — to cram; to stuff (into)}

\label{lesson:GU-C395-thasvu}



\begin{quote}
Before the new one: say the Gujarati for a reservation (of a seat), then the Gujarati for food packed for a journey.
\end{quote}

\subsection*{You'll want to know: \gu{ઠાંસવું}}
\textbf{\gu{ઠાંસવું}} — \emph{ṭhā̃svũ} — \textquotedblleft{}to cram; to stuff (into)\textquotedblright{}.

\begin{quote}
\gu{તેણે} \gu{બધાં} \gu{કપડાં} \gu{એક} \gu{થેલીમાં} \gu{ઠાંસ્યાં}.
\end{quote}

\begin{quote}
\emph{teṇe badhā̃ kapḍā̃ ek thelīmā̃ ṭhā̃syā̃.}
\end{quote}

\begin{quote}
He crammed all the clothes into one bag.
\end{quote}

Past, with \gu{તેણે} (he or she …): \gu{તેણે} \gu{કપડાં} \gu{ઠાંસ્યાં}, \emph{teṇe kapḍā̃ ṭhā̃syā̃}.

\begin{cousinweb}
\textbf{\gu{ઠાંસવું}} (\emph{ṭhā̃svũ}) goes back to \emph{*ṭhass}, \textquotedblleft{}to press down\textquotedblright{}, a form CDIAL reconstructs; CDIAL lists the Gujarati verb, \textquotedblleft{}to cram\textquotedblright{} (CDIAL 5499).
\end{cousinweb}

\begin{grammarlens}[title={What you've built}]
One more word for getting around.
\end{grammarlens}

\subsection*{Your turn}
\begin{itemize}
\raggedright
  \item \emph{Say it:} \emph{ṭhā̃svũ}
  \item \emph{Say it:} \emph{ṭhā̃svũ}, once more
  \item \emph{Say it:} \emph{teṇe badhā̃ kapḍā̃ ek thelīmā̃ ṭhā̃syā̃.} — He crammed all the clothes into one bag.
  \item \emph{Recall:} say the Gujarati for a reservation (of a seat), then the Gujarati for food packed for a journey, then say \emph{ṭhā̃svũ} again
\end{itemize}

\begin{tcolorbox}[breakable,colback=teal!4,colframe=teal!35!black,title={Before you move on}]
What does \gu{ઠાંસવું} mean? (\textquotedblleft{}To cram; to stuff (into)\textquotedblright{}.) Say it once more.
\end{tcolorbox}
\section[\texorpdfstring{\gu{બેસાડવું} (besāḍvũ)}{બેસાડવું (besāḍvũ)}]{\gu{બેસાડવું} (besāḍvũ) — to seat (someone); to make (someone) sit}

\label{lesson:GU-C395-besadvu}



\begin{quote}
Before the new one: say the Gujarati for food packed for a journey, then the Gujarati for to cram.
\end{quote}

\subsection*{You'll want to know: \gu{બેસાડવું}}
\textbf{\gu{બેસાડવું}} — \emph{besāḍvũ} — \textquotedblleft{}to seat (someone); to make (someone) sit\textquotedblright{}.

\begin{quote}
\gu{માતા} \gu{બાળકને} \gu{બારી} \gu{પાસે} \gu{બેસાડે} \gu{છે}.
\end{quote}

\begin{quote}
\emph{mātā bāḷakne bārī pāse besāḍe che.}
\end{quote}

\begin{quote}
The mother seats the child by the window.
\end{quote}

Past, with \gu{તેણે} (he or she …): \gu{તેણે} \gu{બાળકને} \gu{બેસાડ્યું}, \emph{teṇe bāḷakne besāḍyũ}.

\begin{cousinweb}
\textbf{\gu{બેસાડવું}} (\emph{besāḍvũ}) is the causative partner of \gu{બેસવું}, \textquotedblleft{}to sit\textquotedblright{}, which goes back to Sanskrit \emph{upaviśati}, \textquotedblleft{}sits down\textquotedblright{} (CDIAL 2245): to make someone sit.
\end{cousinweb}

\begin{grammarlens}[title={What you've built}]
One more word for getting around.
\end{grammarlens}

\subsection*{Your turn}
\begin{itemize}
\raggedright
  \item \emph{Say it:} \emph{besāḍvũ}
  \item \emph{Say it:} \emph{besāḍvũ}, once more
  \item \emph{Say it:} \emph{mātā bāḷakne bārī pāse besāḍe che.} — The mother seats the child by the window.
  \item \emph{Recall:} say the Gujarati for food packed for a journey, then the Gujarati for to cram, then say \emph{besāḍvũ} again
\end{itemize}

\begin{tcolorbox}[breakable,colback=teal!4,colframe=teal!35!black,title={Before you move on}]
What does \gu{બેસાડવું} mean? (\textquotedblleft{}To seat (someone); to make (someone) sit\textquotedblright{}.) Say it once more.
\end{tcolorbox}
